\section{Thesis Organization}

The remainder of this thesis is organized as follows. Chapter~\ref{sec:relatedwork} surveys related work in computational bias detection, transformer-based text classification, and metalinguistic feature extraction, situating the contributions of this thesis within both the NLP and computational social science literatures. Chapter~\ref{sec:baselines} establishes experimental baselines through human annotation studies and evaluation of existing transformer and large language model approaches; these baselines define the performance ceiling and source-leakage profile against which all subsequent experiments are measured. Chapter~\ref{sec:dataset} describes our extensions to the AllSides dataset, which provide the data scale required for the continued pre-training experiments in Chapter~\ref{sec:continued-pretraining}. Chapter~\ref{sec:continued-pretraining} details our continued pre-training methodology incorporating metalinguistic auxiliary objectives, and motivates the LLM comparison that follows. Chapter~\ref{sec:llm-finetuning} presents comprehensive experiments with fine-tuned large language models and prompt optimization techniques, asking whether general-purpose LLMs can match or exceed specialized continued pre-training. Chapter~\ref{sec:causal-inference} provides a causal analysis of theme-tone-ideology relationships at both topic and community levels using double machine learning and mediation analysis across five annotation paradigms. The central finding is that topic sentiment carries a robust causal relationship with ideology classification on inherently polarising topics, with seven topics replicating across paradigms, but that zero-shot LLMs systematically inflate effect magnitudes while fine-tuning compresses them toward human scale; classification accuracy does not predict causal fidelity. Chapter~\ref{sec:discussion} synthesizes the metalinguistic, comparative, and causal results into a unified picture of where each approach succeeds and fails. Chapter~\ref{sec:limitations-future-work} concludes with limitations and directions for future research.
